% 27/03/2026 : Ajout de Guillaume Lehir
\subsection{IAgraphie}

\subsubsection{Utilisation de l'intelligence artificielle}

L'intelligence artificielle (qui inclut, sans s'y limiter, les grands modèles de langage) peut être utilisée de manière responsable et éthique tout au long d'une étude scientifique. Le principe directeur concernant l'utilisation de l'IA est que les auteurs d'un manuscrit sont entièrement responsables de l'intégrité de leur travail, qu'ils aient utilisé ou non des outils d'IA. 

\subsubsection{Amélioration linguistique}

Les outils d'intelligence artificielle peuvent être utiles pour améliorer la qualité rédactionnelle d'un manuscrit. Il n'est \textbf{pas nécessaire de déclarer} l'utilisation d'outils d'amélioration linguistique basés sur l'IA. 

\subsubsection{Génération d'images}

L'utilisation d'\textbf{images générées par IA est interdite} sauf si elle s'inscrit dans le cadre de traitement de données (cas ci-dessous).

\subsubsection{Outils d'analyse scientifique}

Les outils d'intelligence artificielle peuvent être utilisés dans le cadre d'une analyse scientifique. Cela inclut notamment la génération de code informatique, l'assistance au débogage, la reconnaissance de motifs ou d'autres méthodes d'analyse de données adaptées à une approche par IA.

\textbf{Cette usage doit être déclaré}. L'auteur doit fournir suffisamment de détails pour permettre la
reproductibilité des résultats générés avec l'IA (IA et la version utilisée, instructions utilisées ou
\mbox{\textit{prompts}, ...).} 


